\documentclass[12pt]{article}
\usepackage[margin=0.72in]{geometry}
\usepackage{fontspec}
\setmainfont{Latin Modern Roman}
\usepackage{microtype}
\usepackage{fancyhdr}
\usepackage{titlesec}
\usepackage{enumitem}
\usepackage{needspace}
\usepackage{xcolor}
\usepackage[hidelinks]{hyperref}
\setlength{\parindent}{0pt}
\setlength{\parskip}{0.36em}
\setlength{\headheight}{15pt}
\titleformat{\section}{\large\bfseries\scshape}{}{0pt}{}
\titleformat{\subsection}{\normalsize\bfseries}{}{0pt}{}
\pagestyle{fancy}
\fancyhf{}
\fancyhead[L]{Whately Logic: Week 1 Test}
\fancyhead[R]{The Three Bells}
\fancyfoot[C]{\thepage}
\renewcommand{\headrulewidth}{0.4pt}

\newcommand{\focusbox}[1]{\par\vspace{0.35em}\noindent\begingroup\setlength{\fboxsep}{6pt}\colorbox{gray!12}{\begin{minipage}{\dimexpr\linewidth-2\fboxsep\relax}#1\end{minipage}}\endgroup\par\vspace{0.45em}}
\newcommand{\studentline}{\par\vspace{0.64em}\hrule\vspace{0.68em}}
\newcommand{\shortline}{\par\vspace{0.42em}\hrule\vspace{0.5em}}
\newcommand{\diagnosisblock}{%
\textbf{Conclusion:}\shortline
\textbf{Stated reason:}\shortline
\textbf{Hidden assumption:}\shortline
\textbf{Does the conclusion follow? Explain briefly.}\studentline\studentline
}

\begin{document}
\begin{center}
{\LARGE\bfseries Week 1 Logic Test}\par\vspace{0.16em}
{\large\scshape The Case of the Three Bells}\par\vspace{0.25em}
{A Whately Puzzle on What Follows From What}\par\vspace{0.45em}
\rule{0.82\textwidth}{0.4pt}
\end{center}

\section*{Name: \rule{2.35in}{0.4pt} \hfill Date: \rule{1.25in}{0.4pt}}

\focusbox{\textbf{Whately's test:} Find the conclusion. Find the stated reason. Bring the hidden assumption into the light. Then ask: does the conclusion really follow? Remember also what logic cannot do: logic can test reasoning, but it cannot supply missing facts, wisdom, virtue, observation, or testimony.}

\section*{The Puzzle Story}
At the Academy of Straight Thinking, three bells hang in the old tower.

\begin{itemize}[leftmargin=*]
\item The \textbf{Silver Bell} is supposed to be rung when a student has spoken wisely.
\item The \textbf{Brass Bell} is rung when a student has spoken confidently.
\item The \textbf{Iron Bell} is rung when a student has won loud applause.
\end{itemize}

One morning, the class finds this note pinned to the tower door:

\focusbox{\centering ``Whoever rang the bell last night must be the wisest student in the Academy.''}

Three students are questioned.

\textbf{Clara says:} ``I rang the Iron Bell. Everyone applauded my speech yesterday, so my argument must have been true.''

\textbf{Miles says:} ``I rang the Brass Bell. I spoke more confidently than anyone else, so I must have had the strongest proof.''

\textbf{Jonas says:} ``I rang the Silver Bell. The rules say the Silver Bell is for wisdom, so if I rang it, I must have spoken wisely.''

Then the Head Tutor adds:

\focusbox{``Be careful. Bells may tell us something, but they may not tell us everything. Logic can test whether a conclusion follows, but it cannot discover facts we have not yet learned.''}

\clearpage

\section*{Part I: Diagnose the Three Claims}
For each student, identify the parts of the argument and test whether the conclusion follows.

\Needspace{0.24\textheight}
\subsection*{1. Clara and the Iron Bell}
\textbf{Claim:} ``Everyone applauded my speech yesterday, so my argument must have been true.''
\diagnosisblock

\Needspace{0.24\textheight}
\subsection*{2. Miles and the Brass Bell}
\textbf{Claim:} ``I spoke more confidently than anyone else, so I must have had the strongest proof.''
\diagnosisblock

\Needspace{0.24\textheight}
\subsection*{3. Jonas and the Silver Bell}
\textbf{Claim:} ``The rules say the Silver Bell is for wisdom, so if I rang it, I must have spoken wisely.''
\diagnosisblock

\Needspace{0.45\textheight}
\section*{Part II: Rank the Reasoning}
Answer in complete thoughts.

\begin{enumerate}[leftmargin=*]
\Needspace{0.14\textheight}\item Which student's reasoning is weakest? Why?\studentline\studentline
\Needspace{0.18\textheight}\item Which student's reasoning is strongest or least weak? Why is it still not complete proof?\studentline\studentline\studentline
\Needspace{0.16\textheight}\item What fact or testimony would the class still need before deciding who actually deserved the Silver Bell?\studentline\studentline
\end{enumerate}

\Needspace{0.42\textheight}
\section*{Part III: What Can Logic Decide?}
Mark each item as \textbf{Logic can test this} or \textbf{Logic alone cannot decide this}. Then give a short reason for two of your choices.

\begin{enumerate}[leftmargin=*]
\item Whether applause proves that an argument is true. \hfill \rule{2.3in}{0.4pt}
\item Whether confidence proves that a speaker has strong proof. \hfill \rule{2.3in}{0.4pt}
\item Whether Jonas actually rang the Silver Bell honestly and correctly. \hfill \rule{2.3in}{0.4pt}
\item Whether a conclusion follows from the reason given. \hfill \rule{2.3in}{0.4pt}
\item Whether a student is wise or virtuous in character. \hfill \rule{2.3in}{0.4pt}
\end{enumerate}

\textbf{Explain two choices:}\studentline\studentline\studentline

\Needspace{0.28\textheight}
\section*{Part IV: Repair One Argument}
Choose Clara, Miles, or Jonas. Rewrite the student's argument so that it becomes more careful. You may narrow the conclusion, add needed evidence, or admit what has not yet been proved.

\textbf{Student chosen:} \rule{1.45in}{0.4pt}

\textbf{Repaired argument:}\studentline\studentline\studentline

\Needspace{0.32\textheight}
\section*{Part V: Macbeth Bridge}
Macbeth hears a prophecy. One part of it comes true. Then he begins to imagine kingship and murder.

Complete the comparison:

\focusbox{Macbeth's dangerous reasoning is like the bell puzzle because he is tempted to treat \rule{1.8in}{0.4pt} as if it proves \rule{1.8in}{0.4pt}.}

Now explain your answer in 4--6 sentences. Use the words \textit{conclusion}, \textit{reason}, \textit{hidden assumption}, and \textit{follows}.

\vspace{1.55in}

\Needspace{0.18\textheight}
\section*{Final Whately Claim}
In one sentence, state what this puzzle teaches about what logic can and cannot do.

\studentline\studentline

\end{document}
